\documentclass[11pt]{article}
\usepackage[T1]{fontenc}
\usepackage{lmodern,amsmath,amssymb,amsthm,booktabs,longtable}
\usepackage[margin=1in]{geometry}
\usepackage[colorlinks=true,linkcolor=blue,urlcolor=blue,citecolor=blue]{hyperref}
\usepackage{microtype,xurl}
\newtheorem{theorem}{Theorem}
\newtheorem{lemma}[theorem]{Lemma}
\newtheorem{proposition}[theorem]{Proposition}
\newcommand{\Z}{\mathbb Z}
\newcommand{\Q}{\mathbb Q}
\newcommand{\OK}{\mathcal O_K}
\newcommand{\pos}[1]{\left(#1\right)_+}
\newcommand{\Tn}{\mathcal T}
\newcommand{\eps}{\varepsilon}

\title{A unified Frobenius determinant proof\\for the tenth discriminant\\
\large Version 0.7.0: \(q<570\)}
\author{Research note prepared with Codex}
\date{26 October 2026}
\hypersetup{pdftitle={A unified Frobenius determinant proof for the tenth discriminant, v0.7.0},
pdfauthor={Research note prepared with Codex},
pdfsubject={A self-contained proof of q < 570 and the class-number-one classification}}
\begin{document}
\maketitle

\begin{abstract}
We prove \(q<570\) for the absolute fundamental discriminant of an
imaginary quadratic field of class number one, and complete the finite
classification. One derivative determinant supplies the entire analytic
argument. An auxiliary inert prime proves nonvanishing; Frobenius and
factorial divisibility, compared with complex size, force a smoothed sum
over split and ramified primes to be large. For \(q\ge2000\), one
factorial refinement and a direct primitive-norm count give a uniform
contradiction. For \(570\le q\le2000\), positive averaging and the
quadratic norm Euler product give the sharper comparison, certified on
five bands by exact rational intervals. No earlier version of this paper
is a proof dependency. The external inputs are standard quadratic ideal
theory and reduction of binary quadratic forms, explicit ordinary-prime
estimates, and classical analysis. We use neither Heegner's method nor
a quasi-Riemann hypothesis or quadratic zero-free theorem.
\end{abstract}

\section{Introduction}
For a negative fundamental discriminant \(D\), put \(q=|D|\) and
\(K=\Q(\sqrt D)\). Tao's challenge asks for effective bounds for a
hypothetical tenth imaginary quadratic field of class number one,
without using Heegner's solution \cite{tao}. His subsequent request for
bounds using lighter machinery motivates the route developed here
\cite{taocomment}.

\begin{theorem}\label{main}
If \(h(K)=1\), then \(q<570\).
\end{theorem}

The determinant strategy traces to OpenAI's manuscript on uniform
exclusion of Landau--Siegel zeros \cite{openai}. We use Green's
derivative-matrix variant \cite{green} and the auxiliary-prime
nonvanishing argument of van Dijk \cite{vandijk}. All the determinant
arguments needed here are proved below. The specialized contributions
from the challenge discussion are cited at their points of use.

This version keeps the cutoff and the five interval certificates of
v6.1 \cite{v6}. Its mathematical change is a new uniform treatment of
\(q\ge2000\) using the same determinant estimates as the bounded
interval. This removes the inherited dependence on two earlier papers
and the intermediate boundary at 4000. The result is one proof with
two final estimates, rather than a sequence of improvements that must
be read cumulatively.

\paragraph{The common quantity.}
Class number one gives two competing descriptions of split primes.
Each is represented by the quadratic norm form, which restricts their
number. A determinant argument, on the other hand, forces a weighted
sum of noninert primes to be large. Throughout the paper that sum is
\[
 \Tn(X)=\sum_{p\text{ split or ramified}}\frac{\log p}{p}f(p/X),
\]
where the nonnegative decreasing weight \(f\) is defined by the exact
floor identity in Section~\ref{divisibility}. The same \(\Tn\) appears
in the determinant comparison, the factorial refinements, and both
upper bounds. This common quantity is the organizing device of the proof.

\paragraph{Proof outline.}
Assume \(h(K)=1\) and \(q\ge570\).
\begin{enumerate}
\item The norm obstruction makes every prime \(p<q/4\) inert.
It also identifies the quadratic lattice and its geometry
(Section~\ref{background}).
\item Given a sufficiently large inert prime \(R\), put
\(S=\lfloor3R/64\rfloor\) and \(J=\lceil R/9\rceil\).
Reduction modulo \(R\) proves that a derivative matrix on a lattice
disk has a nonzero minor (Section~\ref{determinant}).
\item Bound this same minor from below by prime divisibility and from
above by complex size. Their comparison gives the common lower bound
for \(\Tn(R)\). An already established bound for \(\Tn\) supplies
additional factorial divisibility, permitting finite refinement
(Sections~\ref{divisibility}--\ref{comparison}).
\item For \(q\ge2000\), a short sieve supplies an inert prime
\(q^3<R\le3q^3/2\). One refinement and a primitive-norm count give
opposing bounds for \(\Tn(R)\). A decreasing scalar function reduces
the entire infinite range to its endpoint \(q=2000\)
(Section~\ref{large}).
\item For \(570\le q\le2000\), a finer auxiliary-prime sieve,
eight finite refinements, and a positive weighted norm-series comparison
give a contradiction on five certified bands
(Sections~\ref{bands}--\ref{certificate}).
\end{enumerate}
The remaining 70 candidates below 570 are handled in
Section~\ref{finite}, yielding the nine values
\(3,4,7,8,11,19,43,67,163\).

\paragraph{Scope of self-containment and notation.}
The rank argument, divisibility transformations, size bounds, norm
estimates, error bounds, and finite finish are included. Standard
quadratic ideal theory and reduced-form theory, the explicitly stated
Rosser--Schoenfeld estimates, and classical summation and special-function
identities are the background inputs. The fixed rational parameter
tables and computational certificates are part of the accompanying
package; they need no earlier version's files. A determinant of enormous
size is used only theoretically: the programs check its numerical
inequalities, not its entries.

All logarithms are natural, \(x_+=\max(x,0)\), and a prime is
\emph{noninert} if it is split or ramified. A \emph{band} is one of
the fixed bounded intervals \([Q_0,Q_1]\) used in the second application.
The common estimates are first formulated for general \(q\ge570\).

\section{Arithmetic and ordinary-prime estimates}\label{background}
\subsection{The norm obstruction and lattice geometry}
Assume henceforth that \(h(K)=1\) and \(q\ge570\).
We first use the classical small-prime obstruction emphasized by
Silberman in the discussion \cite{silberman}.
Every nonrational algebraic integer has norm at least \(q/4\).
A noninert rational prime has an integral ideal of norm \(p\);
under class number one its generator is nonrational. Consequently
\begin{equation}\label{small-inert}
 p<q/4\quad\Longrightarrow\quad p\text{ is inert}.
\end{equation}
In particular \(q\) is an odd prime: an even discriminant has a small
ramified prime 2, and a composite odd discriminant has a ramified prime
at most \(\sqrt q<q/4\). Inertness of 2 gives \(q\equiv3\pmod8\).
The only units are \(\pm1\), and
\[
 \OK=\Z+\Z\omega,\qquad \omega=(1+i\sqrt q)/2,\qquad
 Q_q(a,b)=|a+b\omega|^2=(a+b/2)^2+qb^2/4.
\]
Write
\[
 v=\sqrt q/2,\qquad \rho=\tfrac12\sqrt{1+v^2}.
\]
The lattice has covolume \(v\) and a rectangular tiling of covering
radius \(\rho\). Its disk count \(N(r)\) therefore satisfies
\begin{equation}\label{disk}
 \frac{\pi\max(r-\rho,0)^2}{v}\le N(r)
 \le\frac{\pi(r+\rho)^2}{v}.
\end{equation}
Every principal ideal of norm \(m\) is a rotation and dilation by
\(\sqrt m\) of this lattice.

\subsection{Explicit bounds for ordinary primes}\label{prime-inputs}
The auxiliary-prime sieves compare the logarithmic weight of all rational primes
with the weight of split primes represented by the norm form. We record
the ordinary-prime estimates needed for this comparison and for the
later valuation bounds.
Put \(\vartheta(x)=\sum_{p\le x}\log p\) and
\(A(x)=\sum_{p\le x}\log p/p\). We use Rosser--Schoenfeld
\cite{rs}, with \(c_0=1.01624\):
\begin{align}
 \vartheta(x)&<c_0x,\label{theta-global}\\
 x\left(1-\frac1{2\log x}\right)&<\vartheta(x)
       \quad(x\ge563),\label{theta-lower}\\
 \vartheta(x)&<x\left(1+\frac1{2\log x}\right)\quad(x>1),\label{theta-upper}\\
 A(x)&>\log x-1.334-\frac1{2\log x}\quad(x\ge319).
 \label{mertens}
\end{align}
The last constant follows from their Theorem 6 and the explicit prime
sum at 319: the certificate checks
\(\log319-A(319)+1/(2\log319)<1.334\).

For sharper bounds at our scales, Table I on printed page 90 gives
pairs \((b,\epsilon_b)\) such that
\((1-\epsilon_b)x<\psi(x)<(1+\epsilon_b)x\) for \(x\ge e^b\).
The corollary to their Theorem 14 gives
\(0\le\psi(x)-\vartheta(x)<1.02\sqrt x+3x^{1/3}\).
Thus, uniformly for \(\log x\ge m\ge b\),
\begin{equation}\label{tabletheta}
 1-\epsilon_b-1.02e^{-m/2}-3e^{-2m/3}
 <\frac{\vartheta(x)}x<1+\epsilon_b.
\end{equation}
Let \(c(m),d(m)\) be any certified upper and lower bounds obtained
by taking the best of these inequalities and
\eqref{theta-global}--\eqref{theta-upper}. The table entries used are
included as exact decimals in \texttt{parameters.json}.

We also use \(A(R)<2\log R\) at our scales. This follows without a
further prime estimate: \(\log(R!)\ge RA(R)-\vartheta(R)\), so
\(A(R)\le\log R+c_0<2\log R\).

\section{The common derivative determinant}\label{determinant}
We first construct the determinant conditional on an inert prime of
suitable size. Its existence will be proved separately for each
application; it is not an extra number-theoretic assumption. The
construction and its nonvanishing argument are the same in both ranges.

\subsection{The derivative matrix and its rank}\label{rank}
We use Green's derivative matrix \cite{green}, with interpolation points
in a disk as suggested by Pratt \cite{pratt}. Choosing the auxiliary
prime itself as the monomial width makes the finite-field degree count
particularly simple.
Take an inert prime \(p_0\), and set \(R=p_0\).
Throughout the proof use the fixed constants
\begin{equation}\label{fixed}
 t=\frac3{64},\qquad C=\frac{27}{64},\qquad
 S=\lfloor tR\rfloor,\qquad J=\lceil R/9\rceil,
 \qquad T=\sqrt{CRv/\pi}+\rho.
\end{equation}
Use all \(n\) lattice points \(z_i\) in the disk of radius \(T\).
Then \(n\ge CR\), so \(nJ\ge tR^2\ge RS\).
The matrix has rows \(0\le r<R,0\le s<S\), columns \((i,k)\),
\(0\le k<J\), and entries
\[
 \partial_z^k(z^r w^s)\big|_{(z_i,\bar z_i)}.
\]
These are ordinary derivatives, without division by factorials.

The next lemma is van Dijk's auxiliary-prime argument \cite{vandijk}.
Green also attributes a related unpublished argument to Pescador
\cite{greencomment}; the proof here uses the public finite-field argument.

\begin{lemma}\label{nonzero}
If \(2T<R\) and \(J\le R\), this matrix has full row rank.
\end{lemma}
\begin{proof}
Reduce modulo \(p_0\OK\), a field of \(p_0^2\) elements. The points
are distinct: a nonzero element of \(p_0\OK\) has modulus at least
\(p_0\), whereas two points in the disk differ by less than \(2T\).
Frobenius gives \(\bar z_i=z_i^{p_0}\).
If a nonzero row combination vanished, substitute \(w=z^{p_0}\) in
its polynomial. The exponents \(r+p_0s\) are distinct, so the result
is a nonzero polynomial of degree at most \(RS-1\). For \(k<J\le p_0\),
its ordinary derivatives agree with the specified \(z\)-derivatives
in characteristic \(p_0\). Each node is thus a zero of multiplicity
at least \(J\), since \(k!\ne0\) for \(k<p_0\). This gives at least
\(nJ\ge RS\) zeros, a contradiction. Full rank after reduction
implies full rank in characteristic zero.
\end{proof}

Fix a nonzero \(RS\)-square minor \(\mathcal D\), and put
\(N=RS\) and \(\Delta=|\mathcal D|^2\in\Z_{>0}\).
Different triangular row changes below may be used for different
primes: they leave this same determinant unchanged.

\section{Prime divisibility of the determinant}\label{divisibility}
We now bound the prime divisibility of the fixed positive integer
\(\Delta\). The baseline bound comes from Frobenius at inert primes.
Its leading term is a sum over all primes minus a smoothed sum over
noninert primes. Higher Frobenius factors and factorials add compatible
divisibility to the same minor. The resulting estimates will be combined
with the size bound in Section~\ref{comparison}.

\paragraph{Uniform parameters for the estimates.}
For a chosen determinant take \(R\ge M\ge2^{18}\) and
\(\eps=1/M\). Choose \(e\ge\rho/\sqrt{CRv/\pi}\) and set
\[
 \nu=C(1+2e)^2,\qquad
 \eps<b_0\le(t-\eps)/\nu,\qquad j=1/9+\eps.
\]
The disk count gives \(n\le\nu R\), hence \(b_0\le S/n\).
The applications choose \(e\) uniformly on their whole ranges.
The small geometry conditions used in Section~\ref{arch} will be
stated there and checked in each application. The choice \(M\ge2^{18}\)
ensures that all ordinary-prime and decreasing-error estimates used below
are within their stated domains.

\subsection{The baseline bound and its smoothing weight}
The floor \(\lfloor r/p\rfloor\) counts how many Frobenius factors
can be taken from a row of degree \(r\). Summing this floor over the
rows produces the following weight.
For \(0<u\le1\), put
\begin{equation}\label{fdef}
 f(u)=2au-a(a+1)u^2,\qquad a=\lfloor1/u\rfloor;
 \qquad f(u)=0\ (u>1),\quad f(0)=1.
\end{equation}
This is continuous, nonnegative, nonincreasing, and at most 1. Also
\begin{equation}\label{fidentity}
 f(u)=2u\sum_{m\ge1}(1-mu)_+,\qquad
 \frac2{R^2}\sum_{r=0}^{R-1}\left\lfloor\frac rp\right\rfloor
       =\frac{f(p/R)}p.
\end{equation}
Define the noninert sum
\begin{equation}\label{Tdef}
 \Tn(X)=\sum_{p\text{ noninert}}\frac{\log p}{p}f(p/X).
\end{equation}
It is nonnegative and nondecreasing in \(X\), and is at most the
ordinary cumulative noninert sum.

Green's triangular change \cite{green}, using
\(a_{r,s}=\min(\lfloor r/p\rfloor,S-1-s)\), replaces \(z^rw^s\) by
\(z^{r-pa_{r,s}}(z^p-w)^{a_{r,s}}w^s\). For inert \(p\), every
ordinary derivative of \(z^p-w\) is divisible by \(p\), so
\[
 v_p(\Delta)\ge2\sum_{r,s}a_{r,s}.
\]
Write this normalized valuation as \(e(R,S,p)\). Counting layers gives
\begin{equation}\label{elayers}
 e(R,S,p)=\frac2{R^2S}\sum_{h=1}^{S-1}(S-h)(R-ph)_+.
\end{equation}
For \(p(S-1)\le R-1\), this equals
\((S-1)/R-p(S^2-1)/(3R^2)\). In general,
\begin{equation}\label{etail}
 e(R,S,p)\ge\frac{f(p/R)}p
       -\frac{R-1}{2RSp}
       -\frac{(R-1)(2R-1)}{6RSp^2}.
\end{equation}
Indeed, sum
\(\sum_{s<S}\min(k,S-1-s)\ge Sk-k(k+1)/2\), then use
\(k(k+1)\le r^2/p^2+r/p\) and \eqref{fidentity}.

The main term in the inert-prime sum in \eqref{etail} is exactly
\[
 \sum_p\frac{\log p}{p}f(p/R)-\Tn(R).
\]
Thus the prime sum that must be bounded from below is \(\Tn(R)\).
Its specific weight will also allow the factorial gains and the final
norm-series comparison to be expressed in terms of \(\Tn\).

\subsubsection{The all-prime term}
We first compare the sum over all primes with \(A(R)\), retaining a
nonnegative correction from the floor remainder.
If \(R=ap+b\), the exact floor identity gives
\[
 \frac{f(p/R)}p=\frac1p-\frac1R+\frac1R\delta_a(p/R),\qquad
 \delta_a(u)=2a+1-u^{-1}-a(a+1)u.
\]
On \([1/(a+1),1/a]\), \(\delta_a\) is nonnegative, vanishes at
the endpoints, and has one maximum. Partial summation with
\(d_*x\le\vartheta(x)\le c_*x\), valid for \(x\ge M/33\), gives
\begin{equation}\label{Gfloor}
 G_{\rm floor}=\max(0,Bd_*-Ac_*),\quad
 A=\frac{\log33-H_{33}+1}{2},\quad
 B=\frac{H_{32}-\log33}{2}.
\end{equation}
Here \(H_n=\sum_{j\le n}1/j\). The increasing and decreasing pieces
contribute integrals
\(\frac12\log(1+1/a)-\frac1{2(a+1)}\) and
\(\frac1{2a}-\frac12\log(1+1/a)\), respectively; sum for
\(1\le a\le32\). Thus
\begin{equation}\label{allf}
 \sum_p\frac{\log p}{p}f(p/R)
 \ge A(R)-c(\log M)+G_{\rm floor}\qquad(R\ge M).
\end{equation}

\subsubsection{The small primes and the quadratic tail}
The loss involving \(p^{-2}\) in \eqref{etail} is needlessly large
for small primes. We evaluate their contribution directly, then remove
the known split and ramified primes from the remaining inert-prime tail.
Put \(H=\lfloor1/t\rfloor=21\),
\[
 a_H=\sum_{p\le H}\frac{\log p}p,\quad
 \theta_H=\sum_{p\le H}\log p,\quad
 W_H=\sum_{p\le H}\log p\,(t-pt^2/3).
\]
Let \(B_H\) be a certified upper bound for
\(\sum_{p>H}\log p/p^2\). Specifically, with \(P=10000\),
partial summation and \eqref{theta-global} give
\[
 B_H=\sum_{H<p\le P}\frac{\log p}{p^2}
               +\frac{2c_0}{P}-\frac{\vartheta(P)}{P^2}
 >\sum_{p>H}\frac{\log p}{p^2}.
\]
Choose any nonnegative certified lower bound
\begin{equation}\label{Ktwo}
 K_2\le\sum_{\substack{p>H\\p\text{ noninert}}}\frac{\log p}{p^2},
 \qquad K_2<B_H.
\end{equation}
Then the inert quadratic tail is at most \(B_H-K_2\).
The choice \(K_2=0\) is always available. The bounded-band application
will improve it using explicitly represented prime norms. Combining \eqref{etail}, \eqref{allf}, the
exact small-prime formula, and \(A(R)<2\log R\), gives
\begin{align}\label{basevaluation}
 \frac{\log\Delta}{NR}\ge{}&
 A(R)-c(\log M)+G_{\rm floor}-\Tn(R)-a_H+W_H\\
 &-\frac{B_H-K_2}{3(t-1/M)}
 -\frac{2\theta_H}{M}-\frac{\log M}{(t-1/M)M}.
 \notag
\end{align}
The next two subsections add further divisibility while preserving this
bound for the same determinant.

\subsection{Higher Frobenius factors}\label{frobenius}
Kerr's suggestion to retain higher Frobenius relations, including the
ninth-power relation at 3, motivates the extra factors here
\cite{kerrcomment}. To add their gains, we must put them into one
triangular row transformation and account for differentiation of the
resulting product.
At an inert prime use
\[
 G=z^p-w,\qquad H_1=z^{p^2}-z,\qquad H_2=w^{p^2}-w.
\]
For each row let
\[
 a=\min(\lfloor r/p\rfloor,S-1-s),\quad
 c=\lfloor(r-pa)/p^2\rfloor,\quad b=\lfloor s/p^2\rfloor,
\]
and replace the row polynomial by
\begin{equation}\label{threefactor}
 z^{r-pa-p^2c}G^aH_1^c\,w^{s-p^2b}H_2^b.
\end{equation}
Its leading term is \(z^rw^s\); all other terms have smaller
lexicographic index and remain in the rectangle, since their
\(z\)-degree is at most \(r\) and their \(w\)-degree at most
\(s+a\le S-1\).

The \(k\)th Hasse derivative in \(z\) is the coefficient of \(X^k\)
in \(F(z+X,w)\); ordinary differentiation multiplies it by \(k!\).
All three factors have values divisible by \(p\). A Hasse derivative
can lose a factor \(G\) only at order \(p\), and a factor \(H_1\)
only at order 1 or \(p^2\). At total order \(k\), the number of lost
factors is therefore at most
\[
 L=\min(k,c)+\left\lfloor\frac{k-\min(k,c)}p\right\rfloor.
\]
Ordinary differentiation supplies \(k!\). The elementary inequality
\[
 c+\lfloor k/p\rfloor-L\ge\lfloor c/p\rfloor
\]
holds both for \(k\ge c\) and \(k<c\). Consequently the entry
valuation is at least
\begin{equation}\label{newval}
 a+b+\lfloor c/p\rfloor+\sum_{e\ge2}\lfloor k/p^e\rfloor.
\end{equation}
The additional residual row factor is
\(\lfloor(r-pa)/p^3\rfloor\). Moreover, beyond this bound one may
retain
\begin{equation}\label{coupled-entry}
 \max\left(0,(1-1/p)(c-k)-1\right).
\end{equation}
For \(k<c\), this follows by bounding \(\lfloor k/p\rfloor\)
below by \(k/p-1\); for \(k\ge c\), \eqref{newval} suffices.

\subsubsection{Averaging the row and column gains}
The minor contains every row but only a selection of the available
columns. Each derivative order occurs in at most \(n\) columns.
This restriction bounds the mean derivative order in any chosen minor;
the parameter \(b_0\) is the uniform lower bound for \(S/n\)
specified above.
Suppose \(R\ge M\), put \(\varepsilon=1/M\), and choose
\(b_0\le S/n\). At most \(n\) columns have each derivative order.
Their mean order is at least \((b_0R-1)/2\), and at most \(j_*R\),
where
\begin{equation}\label{jstar}
 j=1/9+\varepsilon,\qquad j_*=j-b_0/2+b_0/M.
\end{equation}
The latter is a conservative consequence of sorting the orders:
the \(i\)th is at most \(J-1-\lfloor(N-1-i)/n\rfloor\).

Let \(P_*\le q/4\); all primes \(p<P_*\) are inert.
The \(b\) and higher-column terms in \eqref{newval} give
\begin{equation}\label{Gsmall}
 G_{\rm small}=\sum_{p<P_*}\log p\left\{
 \sum_{\substack{e\ge2\\p^e<M}}
       \left(\frac{b_0-\varepsilon}{p^e}-2\varepsilon\right)_+
 +\left(\frac{t-\varepsilon}{p^2}-3\varepsilon\right)_+
 \right\}.
\end{equation}
For the residual factor, define
\[
 A_p(t)=
 \begin{cases}
  1/(2p^2)-t/(2p)+t^2/6,&pt\le1,\\
  1/(6tp^3),&pt\ge1.
 \end{cases}
\]
Then the further gains are
\begin{align}
 G_{\rm res}&=\sum_{p<P_*}\log p
       \left(\frac{2A_p(t)}p-(2+2/p^3)\varepsilon\right)_+,
       \label{Gres}\\
 G_{\rm coup}&=2\sum_{p<P_*}\log p
       \left((1-1/p)\{A_p(t)-(1+1/p^2)\varepsilon-j_*\}
                         -\varepsilon\right)_+.
       \label{Gcoup}
\end{align}
Here is a direct integral justification. Put \(x_r=r/R\) and
\(y_s=(S-1-s)/R\). We have
\(c/R\ge (x_r-py_s)_+/p^2-1/R\).
On the cell \([x_r,x_r+1/R]\times[y_s,y_s+1/R]\),
\((x-py)_+\le(x_r-py_s)_++1/R\).
The average of \((x-py)_+/p^2\) over
\([0,1]\times[0,S/R]\) equals \(A_p(S/R)\), which decreases
as \(S/R\) increases. Thus the normalized mean of
\(c\) is at least \(A_p(t)-(1+1/p^2)/M\).
For \eqref{Gres}, use the exact identity
\(\lfloor c/p\rfloor=\lfloor(r-pa)/p^3\rfloor\).
The same cell argument, now divided by \(p^3\), bounds its
normalized mean below by \(A_p(t)/p-(1+1/p^3)/M\).
The mean of a positive part is at least the positive
part of its mean, so \eqref{coupled-entry} and \eqref{jstar} give
\eqref{Gcoup}. All these terms are added to \eqref{basevaluation}.

\subsection{Factorial divisibility and finite refinement}\label{factorials}
Noninert primes contribute factorial divisibility of their own. The
floor identity \eqref{fidentity} expresses the first-power contribution
using the same function \(\Tn\), now at smaller scales. Thus an already
proved lower bound for \(\Tn\) strengthens the next determinant comparison.

At a split prime, replace \(w^s\) by its falling factorial \((w)_s\).
The two components of
\(\OK/p^e\OK\simeq(\Z/p^e\Z)^2\) show that its values are
divisible by \(p^{v_p(s!)}\): in each component a product of \(s\)
consecutive integers has that divisibility. The column factor \(k!\)
remains. At the ramified prime retain only the column factorial.
These are triangular row changes, so they bound the same determinant
as the inert-prime transformations, at a disjoint set of primes.

The split row factorial contributes exactly
\((S/R)\Tn_{\rm split}(S)\) at the first powers of the primes.
For the columns put \(K_0=\lfloor b_0R\rfloor\). The increasing
chosen orders are at least \(\lfloor i/n\rfloor\). Since
\(N/n\ge K_0\), their mean \(\lfloor k/p\rfloor\) is at least
the mean over \(0\le k<K_0\). The identity \eqref{fidentity}
therefore gives a contribution at least
\((K_0/R)\Tn(K_0)\).

Consequently, for any
\(0<z\le b_0-\eps\) and \(0<h\le t-\eps\), a compatible
additional normalized valuation is at least
\begin{equation}\label{Gpointwise}
 z\Tn(zR)+h\left(\Tn(hR)-\frac{\log q}{q}\right)_+.
\end{equation}
Here \(\eps=1/M\) and \(R\ge M\); the ramified term was bounded
by \((\log q)/q\), and \(\Tn\) is nonnegative and nondecreasing.
This pointwise version is useful in the infinite range.

For an interval of scales, suppose a nonnegative nondecreasing function
\(\Phi\) is already known to satisfy \(\Tn(e^x)\ge\Phi(x)\)
at the scales used. Taking \(z=b_0-\eps\), \(h=t-\eps\), and
\(R\ge M\) in \eqref{Gpointwise} gives the uniform gain
\begin{align}\label{Gfeedback}
 G_{\rm fb}={}&(b_0-\eps)\Phi(\log(M(b_0-\eps)))\\
 &+(t-\eps)\left(\Phi(\log(M(t-\eps)))-\frac{\log q}{q}\right)_+.
 \notag
\end{align}
On a band, \((\log q)/q\le(\log Q_0)/Q_0\), giving the constant
used by its certificate. Starting with \(\Phi=0\) makes every later
refinement an implication from a previously proved bound. The infinite
range below needs one refinement; the bounded calculation uses eight.

\section{An upper bound for the complex size}\label{arch}
We now bound the same nonzero minor from above. A typical entry contains
both a falling factorial from differentiation and a power of the radius
of its interpolation point. These depend on the same monomial exponent.
The following estimate keeps that dependence when averaging over rows
and columns. The result is \eqref{arch-final}; it will be compared with
the divisibility bounds in Section~\ref{comparison}.

\subsection{A pointwise inequality for each entry}
Fix \(\tau=3/10\), and, for \(0\le k\le u\le1\), put
\[
 F(u,k)=u\log u-(u-k)\log(u-k)-k,\qquad
 \phi(u)=\tau(u\log u-u).
\]
At zero use the continuous values. For \(0<a\le1\), define
\(\psi(k,a)=\max_{k\le u\le1}\{(u-k)\log a+F(u,k)-\phi(u)\}\).
The maximized function is strictly concave. If \(a\le1-k\), its
critical point satisfies \(u-k=au^{1-\tau}\), and
\[
 \psi(k,a)=\tau u+(1-\tau)k\log u-k.
\]
If \(a\ge1-k\), the maximum is at 1 and
\[
 \psi(k,a)=(1-k)\log\frac a{1-k}-k+\tau.
\]
Differentiation shows that \(\psi\) decreases in \(k\), increases in
\(a\), and is 1-Lipschitz in \(a\). In particular, for \(k\le j\),
\[
 \psi(k,a)+B\ge0,\qquad B=(1-\tau)j(1-\log j).
\]

To prepare for averaging over the disk and derivative orders, let
\(I(k)=\int_0^1 2a\psi(k,a)\,da\) and
\(K(j)=\int_0^j I(k)\,dk\). We have \(I(0)=2\tau^2/(2\tau+1)\).
Changing variables from \(a\) to the critical point gives
\[
 I'(k)=-(1-\tau)\int_k^1(u-k)^2u^{2\tau-3}\,du
               +\log(1-k)+k-k^2/2.
\]
Thus, writing \(A=2\tau\),
\begin{align}\label{Kformula}
 K(j)={}&j\frac{2\tau^2}{2\tau+1}
 -(1-\tau)\left\{\frac{j^2}{2A}-\frac{j^3}{3(A-1)}
 +\frac{j^4}{12(A-2)}
 -\frac{2j^{A+2}}{A(A-1)(A-2)(A+1)(A+2)}\right\}\\
 &+\frac j2-\frac{3j^2}4+\frac{(1-j)^2}2\log(1-j)
        +\frac{j^3}6-\frac{j^4}{24}.
 \notag
\end{align}

Write \((r)_k=r(r-1)\cdots(r-k+1)\) for the falling factorial,
with \((r)_0=1\). For an actual row and derivative order, take \(u=(r+1)/R\) and
\(k'=k/R\). Increasing-log integration gives
\(\log(r)_k\le k\log R+RF(u,k')\).
Clip the normalized radius to \(a=\max(1,|z|)/T\). A nonzero entry
satisfies
\begin{align*}
 \log|\partial_z^k(z^rw^s)|\le{}&
 (r+s)\log T+k\log(R/T)\\
 &+R\{\phi(u)+\psi(k',a)\}+s\log a-\log a.
\end{align*}
The final term is at most \(\log T\). Since \(\phi\) decreases, its
right Riemann sum gives a normalized row contribution at most
\(-3\tau/2\) to \(\log\Delta/(NR)\).

\subsection{Averaging over the lattice and chosen columns}\label{arch-averaging}
The minor may select any \(N\) of the available columns. We therefore
need an upper bound that is uniform over that selection. Adding \(B\)
makes the quantity to be summed nonnegative, so we can first sum over
all available columns and then subtract \(NB\).

With the uniform parameters of Section~\ref{divisibility}, put
\begin{equation}\label{geometry}
 r_* =\sqrt{CRv/\pi},\qquad e\ge\rho/r_*,\qquad
 \nu=C(1+2e)^2,\qquad
 g=\tfrac12\log(C/(2\pi))+\log(1+e).
\end{equation}
Then \(n\le\nu R\), \(b_0\le S/n\), and
\(\log T\le\frac12\log R+\frac14\log q+g\).
These estimates require only the displayed inequality for the actual
\(q,R\). They impose no upper bound on \(q\).

Lattice-cell integration, using nonnegativity and the Lipschitz property,
gives for each order
\[
 \sum_{|z|\le T}\{\psi(k/R,a_z)+B\}
 \le\nu R\left\{I(k/R)+B+\frac{5e}{3(1+e)}\right\}.
\]
Indeed, cells lie in the disk of radius \(T+\rho\). Its normalized
radial argument is \(((1+2e)a+e)/(1+e)\), whose difference from
\(a\) is \(e(1+a)/(1+e)\). Integrating against \(2a\,da\) gives
the displayed error. Clipping at 1 only improves the estimate;
\(\rho\ge1\) also covers the clipped origin.

We may fill in unchosen columns before subtracting \(NB\).
Since \(I\) decreases and \(I+B\ge0\), the order sum is at most
\[
 R\left\{K(j)+j\left(B+\frac{5e}{3(1+e)}\right)\right\}+I(0)+B.
\]
It follows that the normalized joint contribution is at most
\begin{equation}\label{Djoint}
 D_*=-\frac{3\tau}2+
 \frac{2\nu}{t-\varepsilon}
 \left\{K(j)+j\left(B+\frac{5e}{3(1+e)}\right)
                  +\varepsilon(I(0)+B)\right\}-2B.
\end{equation}

\subsection{The remaining radius factor and the final bound}
The entry estimate also contains \(s\log a\), which has not yet been
used. It is nonpositive. To retain a definite saving, count how many
chosen columns must lie inside each smaller disk and how small their
assigned row exponents can be in a determinant term.
At least \(\lambda(a)N\) chosen columns have radius at most \(aT\), where
\[
 \lambda(a)=\left(1-\frac{jC}{t-\varepsilon}
   \{(1+2e)^2-\max(a(1+e)-e,0)^2\}\right)_+.
\]
Among that many distinct rows the normalized sum of the \(s\)'s is
at least \(((t-\varepsilon)\lambda(a)^2-\varepsilon)/2\).
Put
\[
 \delta=\frac{jC}{t-\varepsilon}
       \left\{(1+2e)^2+\frac e{2+e}\right\}-1.
\]
Since \(\max(a(1+e)-e,0)^2\ge a^2-e/(2+e)\) and
\(jC/(t-\varepsilon)\ge1\), we have
\(\lambda(a)\ge(a^2-\delta)_+\).
Assume \(0<\delta<1\) and \(\delta\ge e\), as checked in both applications. Then
\(\sqrt\delta\ge1/T\). Integrating from \(\sqrt\delta\) to 1 gives
the additional saving
\begin{equation}\label{Gs}
 G_s=\left((t-\varepsilon)
       \left\{\frac14-\delta+\frac{3\delta^2}4
                    -\frac{\delta^2}2\log\delta\right\}
                     +\frac\varepsilon2\log\delta\right)_+.
\end{equation}

There are at most \(N!\) terms in the determinant. Using
\(\log N!\le N\log N\), \eqref{jstar}, and \(1<T<R\), we obtain
\begin{equation}\label{arch-final}
 \frac{\log\Delta}{NR}\le
 (1-\alpha)\log R+\frac\beta4\log q+\beta g
       +D_*+\frac{2\log M}{M}-G_s+\frac{2\log(tM^2)}M,
\end{equation}
where
\(\alpha=(1-t)/2-j_*\), \(\beta=1+t-2j_*\).
The functions \((\log R)/R\) and \(\log(tR^2)/R\) decrease for \(R\ge M\ge2^{18}\). The terms \(D_*\) and \(-G_s\) account for different
parts of the displayed entry bound: the former treats
\(\phi+\psi\), and the latter treats \(s\log a\).

\section{The common prime-sum comparison}\label{comparison}
We can now state the conclusion of the determinant argument without
committing to a particular range of discriminants. The main valuation
is an ordinary-prime contribution minus \(\Tn(R)\); comparing it
with the size bound therefore forces \(\Tn(R)\) to be large.

\begin{proposition}\label{common-proposition}
Suppose \(q\ge570\), \(h(K)=1\), and \(R\) is inert. Use
\eqref{fixed} and the uniform parameters specified above, with
\(2T<R\), \(1<T\), \(0<\delta<1\), \(\delta\ge e\),
and \(0<\alpha,\beta\). Choose \(P_*\le q/4\) and \(K_2\)
as in \eqref{Ktwo}. Define
\begin{equation}\label{basecost}
\begin{aligned}
 \mathcal K_0={}&a_H-W_H+c(\log M)
 +\frac{B_H-K_2}{3(t-\eps)}+\beta g+D_*
 +\frac{2\log M}{M}-G_s\\
 &+\frac{2\theta_H}{M}+\frac{\log M}{(t-\eps)M}
 +\frac{2\log(tM^2)}M+1.334+\frac1{2\log M}\\
 &-G_{\rm floor}-G_{\rm small}-G_{\rm res}-G_{\rm coup}.
\end{aligned}
\end{equation}
Then
\begin{equation}\label{common-line}
 \Tn(R)\ge\alpha\log R-\frac\beta4\log q-\mathcal K_0+G,
\end{equation}
where \(G=0\) is always allowed. Any previously established lower
bound for \eqref{Gpointwise}, or \(G_{\rm fb}\) in
\eqref{Gfeedback}, is also allowed as \(G\).
\end{proposition}
\begin{proof}
Add the compatible gains \eqref{Gsmall}, \eqref{Gres}, and
\eqref{Gcoup} to \eqref{basevaluation}, use \eqref{mertens}, and
compare with \eqref{arch-final}. All these bounds concern the same
positive integer \(\Delta\). This gives \eqref{common-line} with
\(G=0\). The additional factorial divisibility in
Section~\ref{factorials} gives the stated refinement.
\end{proof}

The estimates used in this proposition are proved above. The remaining
work consists of supplying inert primes, choosing the uniform constants,
and bounding this same \(\Tn\) from above. These tasks differ between
the infinite range and the interval close to the cutoff.

\section{The uniform exclusion for \texorpdfstring{\(q\ge2000\)}{q >= 2000}}\label{large}
The common proposition already has enough strength in the infinite
range. We use a simple auxiliary-prime interval and one factorial
refinement, then compare the resulting pointwise bound with an
elementary primitive-norm count. This section is independent of the
bounded-band certificates that follow.

Assume \(q\ge2000\), and write \(L=\log q\), \(u=\sqrt q\),
and \(Q=Q_q\) for the norm form.

\subsection{A supply of auxiliary inert primes}

Set $\kappa=3/2$. We claim that, for every
\begin{equation}\label{Urange}
 q^2\le U\le\kappa q^3,
\end{equation}
there is an inert prime $R\in(U,\kappa U]$.

Following Pratt's norm-counting suggestion \cite{pratt}, we use
the elementary condition \(\gcd(a,b,210)=1\).
The disk estimate \eqref{disk} bounds the absolute counting error by
\((2\pi\rho r+\pi\rho^2)/v\), even when \(r<\rho\).
Apply it to both annulus boundaries at each divisor of 210 and use
\[
 \sum_{d\mid210}\frac{\mu(d)}{d^2}=\frac{768}{1225},\quad
 \sum_{d\mid210}\frac{|\mu(d)|}{d}=\frac{96}{35},\quad
 \sum_{d\mid210}|\mu(d)|=16.
\]
Here \(\mu\) is the M\"obius function; for a divisor of 210 its value
is \((-1)^k\), where \(k\) is the number of prime factors.
Every split prime norm has four generators. Consequently the number
of split primes in $(U,\kappa U]$ is at most
\[
 A_q(\kappa-1)U+B_q^*(1+\sqrt\kappa)\sqrt U+C_q^*,
\]
where
\[
 A_q=\frac{\pi}{2u}\frac{768}{1225},\quad
 B_q^*=\frac\pi4\sqrt{1+4/q}\frac{96}{35},\quad
 C_q^*=2\pi(v+v^{-1}).
\]
The area term is exact; using the disk error on both annulus
boundaries gives the displayed square-root and constant errors.

Let $q_0=2000$, $M=q_0^2=4\,000\,000$, and $c_0=1.01624$.
Rosser--Schoenfeld's ordinary-prime estimates give
\[
 \vartheta(\kappa U)-\vartheta(U)
 >\left[\kappa\left(1-\frac1{2\log M}\right)-c_0\right]U.
\]
In contrast the split-prime weight, divided by $U$, is at most
\begin{equation}\label{sievecost}
 (3L+2\log\kappa)
 \left\{A_q(\kappa-1)+\frac{B_q^*(1+\sqrt\kappa)}q
                         +\frac{C_q^*}{q^2}\right\}.
\end{equation}
Each term of this expression decreases for $q\ge2000$.
For the powers of $q$ this follows from
$3-a(3L+2\log\kappa)<0$ for $a\ge1/2$; the extra factor
$\sqrt{1+4/q}$ also decreases. At $q=2000$ the difference between
the ordinary-prime lower bound and \eqref{sievecost} exceeds $0.116$.
The unique ramified prime $q$ is below the annulus. An inert prime
therefore exists throughout \eqref{Urange}.

\subsection{Uniform determinant constants}
Take any inert prime \(R\ge q^2\). Use the dimensions and disk in
\eqref{fixed}, and set
\[
 M=4\,000\,000,\quad \eps=1/M,\quad
 e=\frac{33}{10000},\quad b_0=\frac{2193}{20000},\quad
 P_*=500,\quad K_2=0.
\]
These constants satisfy the common proposition for all \(q\ge2000\).
For the geometric condition,
\[
 \frac{\rho}{\sqrt{CRv/\pi}}
 \le\sqrt{\frac{\pi}{4Cq^2}(v+v^{-1})}<e.
\]
The square of the middle expression decreases with \(q\), so the
endpoint 2000 suffices. The same endpoint verifies \(2T<R\) and
\(T>1\), and
\(b_0\le(t-\eps)/(C(1+2e)^2)\). Notice that this uses
\(R\ge q^2\), not an upper bound on \(q\).

In \eqref{basecost} use the global constant \(c(\log M)=1.01624\).
For \eqref{Gfloor} take
\(d_*=1-1/(2\log(M/33))\) and \(c_*=1.01624\).
The higher-Frobenius sums use only the fixed primes \(p<500\),
all inert here. The exact rational certificate proves
\[
 \alpha>0.42,\qquad 0<\beta/4<0.234,\qquad
 \mathcal K_0<2.211<2.22.
\]
It checks the remaining geometry conditions as well. The actual
enclosure gives \(\mathcal K_0<2.210699\); the larger rounded
constant is convenient for the argument. Proposition~\ref{common-proposition}
with \(G=0\) therefore yields
\begin{equation}\label{base}
 \Tn(R)>0.42\log R-0.234L-2.22
 \qquad(R\text{ inert prime},\ R\ge q^2).
\end{equation}

\subsection{One factorial refinement}

For every real $X$ in
\[
 \kappa q^2\le X\le\kappa^2 q^3,
\]
apply the auxiliary-prime lemma at $U=X/\kappa$. Monotonicity of
$\Tn$ and \eqref{base} imply
\begin{equation}\label{transfer}
 \Tn(X)>0.42\log X-0.234L-2.4,
\end{equation}
because $2.22+0.42\log(3/2)<2.4$.

Now choose an inert prime $q^3<R\le\kappa q^3$, and use its
determinant once more. Put
\[
 z=0.109<b_0-\eps,\qquad h=0.046<t-\eps.
\]
The column factorial contribution is at least $z\Tn(zR)$,
and the split row factorial contribution is at least
$h(\Tn(hR)-L/q)$, by \eqref{Gpointwise}.
Both scales lie in the interval of \eqref{transfer}:
$h q>\kappa$ for $q\ge2000$, and their upper bounds are immediate.
Both lower bounds are positive, and $L/q<0.004$.
Consequently a compatible gain for this same determinant is
\begin{align*}
 G_{\rm fb}\ge{}&z\{0.42\log(zR)-0.234L-2.4\}\\
 &+h\{0.42\log(hR)-0.234L-2.404\}.
\end{align*}
There is no circular step: \eqref{transfer} was established with
zero noninert factorial gain before using this refinement.

Adding this to \eqref{base}, and using $\log R>3L$, gives
\begin{equation}\label{final-lower}
 \boxed{\ \Tn(R)>1.185L-2.76.\ }
\end{equation}
Indeed the exact coefficient is
$(3\cdot0.42-0.234)(1+z+h)=1.18503$, and the constant is
\[
 2.22+2.4(z+h)-0.42(z\log z+h\log h)+0.004h
 <2.753140<2.76.
\]
This uses one refinement, not the eight used near $q=570$.

\subsection{A direct upper bound for the same smoothed sum}

We now give the primitive norm count used in this application.
It implements the split-prime counting approach suggested by Pratt
\cite{pratt}. Write $N(x)$ for the number of pairs
$b>0$, $a\in\mathbb Z$ with $Q(a,b)\le x$.
The horizontal row length is \(2\sqrt{x-qb^2/4}\) when nonempty;
its integral over \(0\le b\le2\sqrt x/u\) is \(\pi x/u\).
The decreasing integral comparison loses at most \(2\sqrt x\),
and the integer-point error at most one per row. Therefore
\[
 \frac{\pi x}{u}-2\sqrt x-\frac{2\sqrt x}{u}
 \le N(x)\le\frac{\pi x}{u}+\frac{2\sqrt x}{u}.
\]
Insert these bounds with the appropriate signs into
$\sum_{d\mid210}\mu(d)N(x/d^2)$. Using the 210 products already
displayed and $\prod_{p\mid210}(1-p^{-1})=8/35$ yields
\[
 N_{210}(x)\le\frac{\pi\eta x}{u}+B_q\sqrt x,
 \qquad \eta=\frac{768}{1225},\qquad
 B_q=\frac{88}{35}+\frac{192}{35u}.
\]
Each split prime has two primitive generators with $b>0$. Hence
\begin{equation}\label{count}
 \pi_+(x)\le\frac{\pi\eta}{2u}x+\frac{B_q}{2}\sqrt x.
\end{equation}
There are no split primes below $q/4$.

There is a useful simple bound for the smoothing weight:
\begin{equation}\label{smooth}
 f(s)\le1-\frac45s\qquad(0\le s\le1).
\end{equation}
For $a=\lfloor1/s\rfloor$, write
$f(s)=1-s+s\delta_a(s)$ with
$\delta_a(s)=2a+1-s^{-1}-a(a+1)s$.
Its maximum is at most
$(\sqrt{a+1}-\sqrt a)^2\le3-2\sqrt2<1/5$, proving
\eqref{smooth}.

For any $X\ge q/4$, the function
$g_X(x)=(\log x)/x\,(1-4x/(5X))$ is nonnegative and decreasing
on $[q/4,X]$. Partial summation with \eqref{count} therefore applies
with an upper bound in the correct direction. Put $x=\log X$ and
$a=\log(q/4)$. The linear part of the count contributes
\[
 \frac{\pi\eta}{4u}
 \left\{x^2-a^2+2a-\frac85\left(x-1+\frac{q}{4X}\right)\right\}.
\]
For the square-root part we may drop the smoothing, giving at most
$2B_q(a+1)/u$; the discarded endpoint term is nonpositive.
The ramified prime contributes at most $L/q$. We obtain
\begin{equation}\label{norm-upper}
 \Tn(X)\le\frac{\pi\eta}{4u}
   \left\{x^2-a^2+2a-\frac85(x-1)\right\}
       +\frac{2B_q(a+1)}u+\frac Lq.
\end{equation}
This bounds exactly the same smoothed noninert sum as
\eqref{final-lower}; no change of target quantity is required.

\subsection{One endpoint covers the infinite range}

Since $R\le\kappa q^3$, monotonicity allows us to use
\eqref{norm-upper} at $X=\kappa q^3$. Write $k=\log(3/2)$ and
$\ell=\log4$. The polynomial in braces, divided by $L$, is
\[
 8L+(6k+2\ell-14/5)
       +\frac{k^2-\ell^2-2\ell-(8/5)k+8/5}{L}.
\]
The rational certificate checks
\[
 \frac{\pi\eta}4<0.4926,\quad
 6k+2\ell-14/5<2.406,\quad
 k^2-\ell^2-2\ell-(8/5)k+8/5<-3.578,\quad
 \ell-1>0.386.
\]
The upper and lower bounds would therefore require
\begin{equation}\label{contradiction}
\begin{aligned}
 1.185<F(q):={}&
 \frac{0.4926}{\sqrt q}\left(8L+2.406-\frac{3.578}{L}\right)
 \\
 &+\frac2{\sqrt q}\left(\frac{88}{35}+\frac{192}{35\sqrt q}\right)
        \left(1-\frac{0.386}{L}\right)
 +\frac1q+\frac{2.76}{L}.
\end{aligned}
\end{equation}
Each term decreases for $L\ge\log2000>7.6$. For the first term,
the sign of the derivative with respect to $L$ is the sign of
\[
 8+\frac{3.578}{L^2}-4L-1.203+\frac{1.789}{L}<0.
\]
For the boundary terms, use
\[
 \frac{0.386}{L^2}-a\left(1-\frac{0.386}{L}\right)<0
 \qquad(a=1/2\text{ or }1).
\]
The remaining terms decrease immediately. Finally the exact interval
calculation proves
\[
 F(2000)<1.166655<1.167<1.185.
\]
This contradicts \eqref{contradiction} for every $q\ge2000$.
The normalized endpoint gap exceeds $0.018$.
The 34 exact checks in \texttt{verify\_tail.py} certify the constants
in this section; the monotonicity arguments above explain why finite
checks cover the infinite range.

\section{Auxiliary primes and refinement on bounded bands}\label{bands}\label{affine}
We now assume \(570\le q\le2000\). Near the cutoff, the direct
count used in Section~\ref{large} is too coarse for the desired bound.
We strengthen the same comparison by working at many scales and then
taking a positive average. All determinant estimates remain those of
Sections~\ref{determinant}--\ref{comparison}.

The sieve in this section implements Pratt's suggestion to use norm
counts to obtain auxiliary inert primes \cite{pratt}. The additional
small-ideal conditions improve the scale intervals available for averaging.
The fixed inputs are precisely those of the retained v6 interval
certificate; we develop their mathematical justification here.

\subsection{Auxiliary inert primes on bounded bands}\label{sieve}
The sieve counts possible split primes in an annulus. Removing elements
divisible by a fixed collection of small ideals improves this count;
a second-order inclusion-exclusion bound keeps the estimate explicit.
Comparing it with the total prime weight will then force an inert prime.

Fix one band \(Q_0\le q\le Q_1\). Form the fixed, finite list of pairs
\begin{equation}\label{normlist}
 b>0,\quad 2a+b>0,\quad \gcd(a,b)=1,\quad
 Q_{Q_1}(a,b)<(Q_0/4)^2,
\end{equation}
ordered by \((Q_{Q_1}(a,b),a,b)\). Write \(p_j(q)=Q_q(a_j,b_j)\),
\(p_j^-=Q_{Q_0}(a_j,b_j)\), and \(p_j^+=Q_{Q_1}(a_j,b_j)\).

Each \(p_j(q)\) is prime. A composite norm below \((q/4)^2\) has a
prime factor below \(q/4\); inertness would force that prime to divide
both coordinates. It is split rather than ramified: generators of the
unique prime ideal above \(q\) have zero real part. These primes are
distinct, because the four generators of a split prime norm are
\(\alpha,-\alpha,\bar\alpha,-\bar\alpha\), and the chosen quadrant
contains exactly one of them.

Choose the first \(n_s\) such split primes, and the first \(h\) rational
primes \(r>7\) with \(r<Q_0/4\). Sieve by the two prime ideals above
each selected \(p_j\), and by each ideal \(r\OK\). Their norms are
\(p_j,p_j,r^2\), and they are pairwise coprime. Also impose the exact
condition \(\gcd(a,b,210)=1\).

If an element belongs to \(k\) of the selected ideals, then
\(1-k+\binom{k}{2}\) is an upper bound for the indicator that it
belongs to none. Apply this second Bonferroni inequality, multiplied
by the exact 210 indicator, to the annulus
\(U<|\alpha|^2\le\kappa U\). A split prime in this interval has four
generators. Assume \(U>Q_1\) and \(U>\max_{j\le n_s}p_j^+\), so the
sieving primes themselves are not counted.

For an ideal of norm \(m\), the disk estimate gives an annulus count
with area term \(\pi(\kappa-1)U/(vm)\), boundary error at most
\(2\pi\rho(1+\sqrt\kappa)\sqrt U/(v\sqrt m)\), and constant error
at most \(2\pi\rho^2/v\). The 210 expansion has
\[
 \sum_{d\mid210}\frac{\mu(d)}{d^2}=\frac{768}{1225},\qquad
 \sum_{d\mid210}\frac{|\mu(d)|}{d}=\frac{96}{35},\qquad
 \sum_{d\mid210}|\mu(d)|=16.
\]
Define
\begin{align*}
 u&=2\sum_{j\le n_s}(p_j^+)^{-1}+\sum_{r}r^{-2},&
 u_2&=2\sum_{j\le n_s}(p_j^+)^{-2}+\sum_r r^{-4},\\
 z&=2\sum_{j\le n_s}(p_j^-)^{-1/2}+\sum_r r^{-1},&
 z_2&=2\sum_{j\le n_s}(p_j^-)^{-1}+\sum_r r^{-2},
 \qquad m_0=2n_s+h.
\end{align*}
The certificate checks that the analogous \(u\) at the minus endpoints
is less than 1. The area polynomial then decreases in every reciprocal
ideal norm, since its partial derivative is \(-1\) plus the sum of
the other reciprocals. Therefore uniform upper bounds for the number
of split primes in \((U,\kappa U]\) are
\begin{equation}\label{split-annulus}
 A_*(\kappa-1)U+B_*(1+\sqrt\kappa)\sqrt U+C_*,
\end{equation}
where \(v_-=\sqrt{Q_0}/2\), \(v_+=\sqrt{Q_1}/2\), and
\begin{align*}
 A_*&=\frac{\pi}{4v_-}\frac{768}{1225}
          \left(1-u+\frac{u^2-u_2}{2}\right),\\
 B_*&=\frac\pi4\sqrt{1+v_-^{-2}}\frac{96}{35}
          \left(1+z+\frac{z^2-z_2}{2}\right),\\
 C_*&=2\pi(v_++v_+^{-1})
          \left(1+\frac{m_0(m_0+1)}2\right).
\end{align*}

For a parameter row, put \(M=2^a\), and let
\(2^a\le U\le2^b\). Set \(d=d(\log M)\), \(c=c(\log M)\).
An inert prime exists in \((U,\kappa U]\) if
\begin{equation}\label{sieve-margin}
 G(x)=\kappa d-c-(x+\log\kappa)
 \{A_*(\kappa-1)+B_*(1+\sqrt\kappa)e^{-x/2}+C_*e^{-x}\}>0,
 \quad x=\log U.
\end{equation}
Indeed, the total prime logarithmic weight exceeds the possible split
weight, and the ramified prime is below the interval. Since
\(x+\log\kappa>4\), the function \(G\) is concave. It suffices to
check its two endpoints. These are exact rational interval checks in
the certificate; all \(\kappa,n_s,h,a,b\) are in the parameter file.

\subsection{The geometry and quadratic-tail subtraction}\label{band-inputs}
For a parameter row with \(M=2^a\), choose
\[
 v_+=\sqrt{Q_1}/2,\qquad \rho_+=\tfrac12\sqrt{1+v_+^2},\qquad
 e\ge\frac{\rho_+}{\sqrt{CMv_+/\pi}},\qquad
 0<b_0\le\frac{t-1/M}{C(1+2e)^2}.
\]
Since \(\rho/\sqrt v=\tfrac12\sqrt{v+v^{-1}}\) increases in
our range, these choices satisfy the common geometric hypothesis for
every \(q\) in the band and \(R\ge M\). Take \(P_*=Q_0/4\).
With all pairs in \eqref{normlist}, use
\begin{equation}\label{band-Ktwo}
 K_2=\sum_j\frac{\log p_j^+}{(p_j^+)^2}
             +\frac{\log Q_1}{Q_1^2}.
\end{equation}
The function \((\log x)/x^2\) decreases at these primes. The prime
norms in the list and the ramified prime therefore prove
\eqref{Ktwo}. The certificate checks \(K_2<B_H\), the auxiliary-prime
margins, and all the remaining hypotheses of the common proposition.

\subsection{From a determinant scale to an interval of scales}
Apply Proposition~\ref{common-proposition}, with the band inputs
just specified and the already proved factorial gain. This gives
\begin{equation}\label{Tline}
 \Tn(R)\ge\alpha\log R-\frac\beta4\log q-\mathcal K,
 \qquad \mathcal K=\mathcal K_0-G_{\rm fb}.
\end{equation}
For a band-uniform calculation replace \((\log q)/q\) in
\eqref{Gfeedback} by \((\log Q_0)/Q_0\).

For every \(U\in[2^a,2^b]\), the auxiliary prime satisfies
\(U<R\le\kappa U\). Monotonicity of \(\Tn\) and \(\alpha>0\)
therefore transfer \eqref{Tline} to \(X=\kappa U\), with the
additional intercept cost \(\alpha\log\kappa\).

The auxiliary prime need not have a fixed ratio to \(U\); the interval
guarantee alone suffices. The certificate rounds slopes down and
intercept costs up, obtaining
\[
 \Tn(e^x)\ge A_i x-D_i-\gamma\log(q/Q_0),\qquad
 \gamma=3/10,
\]
on the indicated finite interval. It checks \(\beta/4<\gamma\).
The term involving \(q\) records the possible loss relative to the
lower band endpoint \(Q_0\).

\subsection{Eight finite refinements}
Each affine bound applies on a finite interval of \(x\). Beyond its
right endpoint, keep its endpoint value: monotonicity of \(\Tn\)
makes that value a valid lower bound at all later scales. Taking the
maximum of the available bounds and zero gives a nondecreasing
envelope.

For each of eight refinement stages, subtract
\(\gamma\log(Q_1/Q_0)\) from the preceding envelope and clip at
zero. This gives a band-uniform lower bound \(\Phi\), which can be
inserted into the factorial gain \eqref{Gfeedback}. Repeating the
determinant comparison produces the next stage. The initial stage
uses \(\Phi=0\); every later stage is an implication from previously
proved bounds, with no limiting or fixed-point assertion.

Write \(L(x)\) for a certified piecewise affine selection from the
final bounds at \(Q_0\), including zero and frozen endpoint values.
Then throughout the band,
\begin{equation}\label{Luniform}
 \Tn(e^x)\ge L(x)-\gamma\log(q/Q_0).
\end{equation}
The exact parameter and output files record all lines, stages, and
selected segments. For the rest of the proof, the conclusion of the
determinant argument is precisely \eqref{Luniform}.

\section{Positive kernels and the quadratic norm series}\label{mellin}
The lower bound \eqref{Luniform} is strongest at different scales in
different bands. We average over the scale using a positive kernel,
then use the norm Euler product to bound the same average from above.
The resulting norm series is estimated by Poisson summation, following
the approach discussed by michaelmross and Tao \cite{ross,taocomment}.

\subsection{A positive average of the determinant lower bound}
Choose positive rational \(d,e\) and a nonzero integer polynomial
\(P\), and normalize
\[
 h(x)=Z^{-1}e^{-dx}P(e^{-ex})^2\quad(x\ge0),\qquad
 Z=\int_0^\infty e^{-dx}P(e^{-ex})^2\,dx.
\]
Then \(h\ge0\) and \(\int h=1\). Expand
\(h(x)=\sum_i c_i d_i e^{-d_i x}\), with \(d_i=d+ie\).
All \(c_i\) and \(Z\) are exact rationals. The coefficients \(c_i\)
may have either sign, but the unexpanded kernel is nonnegative.
Integrating
\eqref{Luniform} gives
\begin{equation}\label{weightedlower}
 \int_0^\infty h(x)\Tn(e^x)\,dx
 \ge\sum_i c_i\!\int_0^\infty d_i e^{-d_i x}L(x)\,dx
             -\gamma\log(q/Q_0).
\end{equation}
One and the same \(L\) is used in every signed term: the inequality
is obtained by integration against \(h\ge0\), before expanding the
square. Termwise lower estimates with unrelated functions \(L\)
would not justify this step.

\subsection{The prime weight and its Euler-product upper bound}
Define
\[
 W(y)=\int_y^\infty h(x)f(e^{y-x})\,dx.
\]
This is nonnegative, nonincreasing, and at most 1. Tonelli's theorem
and \eqref{fidentity} show that
\begin{equation}\label{multiplier}
 W(y)=\sum_i c_i F(d_i)e^{-d_i y},\qquad
 F(d)=\int_0^1 d\,u^{d-1}f(u)\,du
       =\frac{2d\zeta(1+d)}{(1+d)(2+d)}.
\end{equation}
The right side follows by integrating each
\(2d\,u^d(1-mu)_+\) from 0 to \(1/m\).
Thus the left side of \eqref{weightedlower} is
\(\sum_{p\text{ noninert}}(\log p)/p\,W(\log p)\).

We next express an upper bound for this weighted prime sum through the
norm series. For \(s>1\), put
\[
 \mathcal A_q(s)=\sum_{b\ge1}\sum_{a\in\Z}Q_q(a,b)^{-s},\qquad
 B_q(s)=1+\mathcal A_q(s)/\zeta(2s).
\]
Every nonzero ideal is principal and has exactly two generators, differing
by sign. Rational generators contribute \(\zeta(2s)\); a nonrational
ideal has exactly one generator with \(b>0\). Thus
\(\zeta_K(s)=\zeta(2s)+\mathcal A_q(s)\). Dividing out rational
integer factors in the ideal Euler product gives
\[
 B_q(s)=\prod_{p\text{ split}}\frac{1+p^{-s}}{1-p^{-s}}\,(1+q^{-s}).
\]
An inert prime contributes 1. At a split prime the primitive ideals
are positive powers of either of the two prime ideals, contributing
\(1+2\sum_{k\ge1}p^{-ks}\); at the ramified prime the primitive
exponents are 0 and 1. All products and differentiated series converge
absolutely for \(s>1\). Logarithmic differentiation now gives
\[
 -\frac{B_q'(s)}{2B_q(s)}
 =\sum_{p\text{ split}}\sum_{\substack{m\ge1\\m\text{ odd}}}
      \frac{\log p}{p^{ms}}
 +\frac{\log q}{2}\sum_{m\ge1}(-1)^{m-1}q^{-ms}.
\]
Multiply by the signed coefficients \(c_iF(d_i)\) and sum at
\(s_i=1+d_i\). Since \(W\ge0\), the discarded higher odd powers
have nonnegative weights. Adding
\(\frac12\log q\sum_{m\ge1}q^{-m}W(m\log q)\) also covers the
ramified first-power term. It is at most
\begin{equation}\label{ramerror}
 \frac{\log Q_0}{2(Q_0-1)}.
\end{equation}

\subsection{Poisson summation with an explicit error}\label{poisson}
The zero Fourier mode of the norm series has a closed form, and the
remaining modes decay exponentially in \(v\). Both the series and
its derivative need explicit error bounds, because the Euler-product
comparison uses a logarithmic derivative.
Set
\[
 C(s)=\sqrt\pi\frac{\Gamma(s-1/2)}{\Gamma(s)},\quad
 g(s)=\psi(s)-\psi(s-1/2),\quad
 \mathcal F_q(s)=C(s)v^{1-2s}\zeta(2s-1).
\]
Here \(\psi=\Gamma'/\Gamma\) is the digamma function, distinct
from the auxiliary function used in Section~\ref{arch}.
Poisson summation in the real coordinate gives
\(\mathcal F_q\) as the zero-mode contribution. We bound the error
uniformly on the band, allowing for the magnitudes of all signed
kernel coefficients.

Choose any rational \(0<a_0<v_-\). For
\(f_{s,t}(x)=(x^2+t^2)^{-s}\), shift the Fourier contour by
\(y=t-a_0\), where \(t=bv\). Since
\[
 |(x+iy)^2+t^2|\ge x^2+a_0^2,
\]
the nonzero Fourier coefficients are bounded by
\(C(s)a_0^{1-2s}e^{-2\pi n(t-a_0)}\). Summing both signs and
all \(b,n\ge1\), put
\[
 z=e^{-2\pi v_-},\quad y_0=e^{-2\pi(v_--a_0)},\qquad
 E_A=\frac{2C(s)a_0^{1-2s}y_0}{(1-z)(1-y_0)}.
\]
Then \(|\mathcal A_q(s)-\mathcal F_q(s)|\le E_A\) uniformly in
the band.

For the \(s\)-derivative, the complex argument of
\((x+iy)^2+t^2\) has absolute value at most \(\pi/2\), and its
logarithm is bounded in modulus by
\[
 \log(1+x^2/a_0^2)+2\log(2t)+2|\log a_0|+\pi/2.
\]
The integral of the first term against
\((x^2+a_0^2)^{-s}\) is
\(C(s)a_0^{1-2s}g(s)\). Using \(\log b\le b-1\) in the remaining
geometric sum therefore gives
\begin{equation}\label{EN}
 |\mathcal A_q'(s)-\mathcal F_q'(s)|\le E_N,
 \quad E_N=E_A\left\{2\log(2v_+)+2|\log a_0|+g(s)+\frac\pi2
                                      +\frac{2z}{1-z}\right\}.
\end{equation}
The contour shift is justified by decay at infinity for \(s>1\);
the same decay with an extra logarithm justifies differentiation.

\subsection{The resulting upper bound}
Put \(d=s-1\), \(z_2=\zeta(2s)\),
\(c_s=-2\zeta'(2s)/\zeta(2s)\), and
\[
 \kappa_d=g(s)-2\frac{\zeta'(2s-1)}{\zeta(2s-1)}-c_s.
\]
The main upper term is
\begin{equation}\label{Hmain}
 H_q(d)=\frac{\mathcal F_q(s)\{\log(q/4)+\kappa_d\}}
                   {2\{z_2+\mathcal F_q(s)\}}.
\end{equation}
The certificate checks that \(H_q(d)>0\) and decreases with \(q\)
for each exponent used. Since \(z_2>1\), the logarithmic-derivative
error is at most
\begin{equation}\label{Herror}
 E_H(d)=\frac{E_N+(c_s+2H_{Q_0}(d))E_A}{2(1-E_A)}.
\end{equation}
Consequently a uniform upper bound for the weighted noninert sum is
\begin{equation}\label{weightedupper}
 U(q)=\sum_i c_iF(d_i)H_q(d_i)
       +\sum_i|c_i|F(d_i)E_H(d_i)
       +\frac{\log Q_0}{2(Q_0-1)}.
\end{equation}
The main terms retain their signs. The error terms use absolute
coefficient values, so cancellation is not assumed in the error
estimate. Equations \eqref{weightedlower} and \eqref{weightedupper}
now bound the same positive average from opposite sides.

\section{Exact certification and completion of the proof}\label{certificate}
The final step has two numerical obligations on each band: a strict
gap at \(Q_0\), and a derivative comparison ensuring that the gap
persists up to \(Q_1\). All parameters are fixed before checking
either obligation.

\subsection{A contradiction on each whole band}
The certificate uses the five bands
\[
 [570,580],\quad[580,600],\quad[600,700],\quad[700,1050],\quad[1050,2000].
\]
The first uses a polynomial of degree 16 with \(d=3/10,e=1/25\);
the others use degree 12 with \(d=2/5,e=1/20\). The integer
polynomial coefficients, auxiliary-prime rows, and eight-stage
refinement schedule are fixed in \texttt{parameters.json}.

The lower bound \eqref{weightedlower} decreases with slope
\(-\gamma\) as a function of \(Y=\log(q/Q_0)\). To compare the
upper bound with it, write
\(a_i=1/2+d_i\) and
\(r_i=\mathcal F_{Q_0}(s_i)/\zeta(2s_i)\). Expand
\[
 H_q(d_i)=\frac12(\log(Q_0/4)+Y+\kappa_{d_i})
       \sum_{\ell\ge1}(-1)^{\ell+1}r_i^\ell e^{-\ell a_iY}.
\]
The certificate bounds the differentiated geometric tail explicitly.
For each retained \(\ell\), the signed finite sum is evaluated as
a polynomial in \(1-e^{-\ell eY}\), with its coefficients computed
by exact interval arithmetic. This avoids naive interval cancellation.
Cells of width \(1/1000\) cover the entire \(Y\)-interval. On every
cell the derivative of the main sum in \eqref{weightedupper} is
strictly less than \(-\gamma\). The two error bounds in that equation
are uniform constants, so they require no differentiation.

\IfFileExists{numerical_table.tex}{\begin{center}
\small
\begin{tabular}{@{}rrrrr@{}}
\toprule
Band & Lower bound & Upper bound & Gap $>$ & Norm error $<$\\
\midrule
$[570,580]$ & 10.975413 & 10.965024 & 0.010389 & $7.652\times10^{-5}$\\
$[580,600]$ & 9.566392 & 9.544454 & 0.021938 & $1.969\times10^{-10}$\\
$[600,700]$ & 9.755174 & 9.682808 & 0.072367 & $4.675\times10^{-11}$\\
$[700,1050]$ & 9.890348 & 9.507108 & 0.383240 & $6.739\times10^{-14}$\\
$[1050,2000]$ & 12.460477 & 11.013702 & 1.446776 & $3.672\times10^{-22}$\\
\bottomrule
\end{tabular}
\end{center}
All displayed bounds are rounded outward. The upper bound includes
both the norm-series error and the ramified-prime allowance.
The exact certificate endpoints are recorded in the JSON output.
}{}

At \(q=Q_0\), the exact lower endpoint of
\eqref{weightedlower} is strictly greater than the exact upper
endpoint of \eqref{weightedupper}. Since the latter decreases faster
than \(-\gamma\log(q/Q_0)\), \eqref{Luniform} proves a contradiction
throughout the band. The five bands exclude \([570,2000]\).
Section~\ref{large} excludes every \(q\ge2000\), proving
Theorem~\ref{main} without an earlier-version dependency.

\subsection{The remaining finite range}\label{finite}
We include the short finite finish, following the small-prime sieving
suggestion in Pratt's comment \cite{prattsieve}.
For \(q>8\), the norm obstruction at 2 forces \(q\equiv3\pmod8\).
For an odd prime \(p<q/4\), inertness requires \(-q\) to be a nonzero
nonsquare modulo \(p\). The following complete integer computation
tests those necessary conditions on the 70 candidates below 570:
\begin{verbatim}
primes = [p for p in range(3, 110)
          if all(p % d for d in range(2, p))]
squares = {p: {a*a % p for a in range(p)} for p in primes}
survivors = [q for q in range(11, 570, 8)
             if all(4*p >= q or (-q) % p not in squares[p]
                    for p in primes)]
assert survivors == [11, 19, 43, 67, 163]
\end{verbatim}
The negative fundamental discriminants of absolute value at most 8 are
\(-3,-4,-7,-8\), giving the nine claimed possibilities.

For the converse, enumerate reduced positive definite primitive forms
\((a,b,c)\) of discriminant \(-q\). Reduction gives
\(1\le a\le\sqrt{q/3}\), \(-a\le b\le a\), and
\(c=(b^2+q)/(4a)\ge a\). Require \(c\) integral,
\(\gcd(a,b,c)=1\), and \(b\ge0\) when \(|b|=a\) or \(a=c\).
For each of the nine values, this finite enumeration gives exactly one
form: \((1,0,q/4)\) for \(q=4,8\), and
\((1,1,(q+1)/4)\) for the seven odd values.
The standard reduced-form correspondence with ideal classes proves:
\begin{theorem}\label{classification}
A negative fundamental discriminant has class number one exactly when
its absolute value is
\[
 3,4,7,8,11,19,43,67,163.
\]
\end{theorem}
The program \texttt{verify\_finite.py} performs both finite steps.
This computation completes the classification below the analytic cutoff;
no discriminants in \([570,\infty)\) were enumerated to prove that cutoff.

\subsection{Reproducing the certificate}
The numerical proof uses integer and rational intervals throughout.
Logarithms use a convergent positive series after range reduction;
exponentials use a Taylor series with an explicit tail; square roots
use integer square roots. Zeta values and derivatives use
Euler--Maclaurin with the periodic-Bernoulli remainder, while gamma
and digamma use positive-real Stirling expansions with remainder
bounds. These are evaluations in \(\Re s>1\), not zero-location
inputs. Appendix~\ref{remainders} specifies the remainder formulas.
Floating-point searches selected the fixed inputs, but acceptance
depends only on the stated interval inequalities.

Run, without Python's \texttt{-O} option,
\begin{verbatim}
python3 verify_tail.py
python3 verify_v6.py
python3 verify_structures.py
python3 verify_finite.py
\end{verbatim}
The first command checks the 34 constant inequalities for the uniform
tail in Section~\ref{large}. The second certifies the five bounded
bands and writes exact endpoints to JSON. Its filename is retained to
make the unchanged v6 calculation easy to compare with its predecessor.
The third checks finite instances of the algebraic identities and
divisibility statements; these tests do not replace their proofs.
The fourth performs the finite check in Section~\ref{finite}.

The bounded-band run passes 429 grouped checks, including 1,258 cells
for the dependence on \(q\), using 184 auxiliary-prime parameter
rows. Its smallest strict gap exceeds \(0.010389\).
The parameter file and exact output files are part of the accompanying
package. The programs check explicit inequalities; they do not
construct the theoretical determinant.


In addition, the uniform-tail certificate gives the normalized gap
\(1.185-F(2000)>0.018345\). The structural checks pass 101,524 finite
cases. All files needed for these runs are included in this version;
the source manifest records their hashes. Numerical certification is
distinct from checking the general mathematical deductions or formalizing
them in a proof assistant.

\section*{Acknowledgements and AI disclosure}
We thank Tao for organizing the challenge and its participants for
sharing their arguments. The determinant strategy originates in OpenAI's
release \cite{openai}; Green and van Dijk developed the variants used
here \cite{green,vandijk}. Green credits ChatGPT for the mathematical
arguments and describes his role as exposition. The other contributions
are cited where they enter the proof. No priority claim is made for the
refinements or consolidation in this note.

Codex developed the refinements, selected parameters, and prepared the
manuscript and programs through an extended collaboration with the
repository owner. V6 received a reported second-model cross-check and a
separate AI-assisted audit; their scope is documented in the repository.
The bounded certificate is retained here, but the new uniform tail and
the consolidated v7 exposition have not yet received a separate model
cross-check. No human peer review or formal verification of v7 is claimed.
The exact computations supplement the mathematical argument; they do
not certify every deduction in the paper. The cutoff 570 is not asserted
to be optimal for this method.

\appendix
\section{Numerical remainder bounds}\label{remainders}
Here are the remainder formulas used in the certificate; the integer
parameters in this paragraph are numerical evaluation parameters,
independent of the determinant dimensions.
For \(s>1\), Euler--Maclaurin with \(N_0=64\), \(m=18\) gives
\begin{align*}
 \zeta(s)={}&\sum_{n<N_0}n^{-s}+\frac{N_0^{1-s}}{s-1}
      +\frac12N_0^{-s}
      +\sum_{k=1}^{m}\frac{B_{2k}(s)_{2k-1}}{(2k)!}
                      N_0^{1-s-2k}+\mathcal R,\\
 |\mathcal R|\le{}&E=
 \frac{|B_{2m}|(s)_{2m}}{(2m)!(s+2m-1)}N_0^{1-s-2m},\\
 |\mathcal R'|\le{}&E\left\{\log N_0+
       \sum_{j=0}^{2m-1}\frac1{s+j}+\frac1{s+2m-1}\right\}.
\end{align*}
Here \((s)_j\) denotes the rising factorial. These bounds follow
from the periodic-Bernoulli integral and its derivative.
For log gamma and digamma, shift to \(z=x+64\), retain the terms
\(1\le k<20\) in the positive-real Stirling expansions, and use
the first omitted term as an absolute remainder bound
\cite{dlmf}:
\[
 |R_{\log\Gamma}|\le\frac{|B_{40}|}{40\cdot39\,z^{39}},
 \qquad |R_\psi|\le\frac{|B_{40}|}{40z^{40}}.
\]
The recurrence formulas then recover the values at \(x\).

After range reduction to \(1\le y\le2\), logarithms use 80 terms
of \(2\sum_{k\ge0}t^{2k+1}/(2k+1)\), where
\(t=(y-1)/(y+1)\), and tail
\(2t^{161}/(161(1-t^2))\). Exponentials are reduced to
\(0\le x\le1\); the series through degree 80 has positive tail
at most \(82/(81\cdot81!)\). Repeated squaring and reciprocation
undo the reduction. Machin's formula with 60 arctangent terms
encloses \(\pi\). Elementary interval operations round outward
to 65 decimal places; square roots use integer bounds.

For the differentiated geometric series in Section~\ref{certificate},
the tail after
\(L\) terms is at most
\[
 \frac12\sum_i |c_i|F(d_i)r_i^{L+1}
 \left\{\frac1{1-r_i}+a_i E_*
       \left(\frac{L+1}{1-r_i}+\frac{r_i}{(1-r_i)^2}\right)\right\},
 \qquad E_*=\max_i\{\log(Q_1/4)+\kappa_{d_i}\}.
\]
The certificate increases \(L\) until this bound is below
\(10^{-10}\), then includes it in every derivative cell test.

\begin{thebibliography}{99}
\bibitem{tao} T.~Tao, \emph{An experimental AI-math challenge: using
tenth-discriminant bounds to explore arguments ruling out Siegel zeroes},
20 October 2026.
\href{https://terrytao.wordpress.com/2026/10/20/an-experimental-ai-math-challenge-using-tenth-discriminant-bounds-to-explore-arguments-ruling-out-siegel-zeroes/}{Challenge and discussion}.
\bibitem{taocomment} T.~Tao, comment 698471 on \cite{tao}.
\href{https://terrytao.wordpress.com/2026/10/20/an-experimental-ai-math-challenge-using-tenth-discriminant-bounds-to-explore-arguments-ruling-out-siegel-zeroes/comment-page-1/\#comment-698471}{Request for bounds using lighter machinery}.
\bibitem{openai} OpenAI, \emph{Uniform exclusion of Landau--Siegel zeros},
manuscript dated 1 October 2026, publicly released 6 October 2026.
\href{https://github.com/openai/math/blob/main/preprints/Uniform-exclusion-of-Landau-Siegel-zeros-October-1-2026/paper.pdf}{Original manuscript};
\href{https://openai.com/index/sharing-ai-progress-in-mathematics/}{release announcement}.
\bibitem{green} B.~Green, \emph{On Landau--Siegel zeros}, v2, 2026.
\href{https://hexagonmath.org/2610.00311v2}{hexagon:2610.00311v2}.
\bibitem{vandijk} R.~van Dijk, \emph{A Frobenius determinant proof of
the non-existence of Landau--Siegel zeros}, v1, 2026.
\href{https://hexagonmath.org/2610.00409v1}{hexagon:2610.00409v1}.
\bibitem{v6} \emph{A smoothed Frobenius determinant bound for the tenth
discriminant}, version 0.6.1, 21 October 2026.
\href{https://github.com/twinheterodox/tenth-discriminant/releases/tag/v0.6.1}{Previous exposition and fixed interval certificate}.
\bibitem{silberman} L.~Silberman, comment 698449 on \cite{tao}.
\href{https://terrytao.wordpress.com/2026/10/20/an-experimental-ai-math-challenge-using-tenth-discriminant-bounds-to-explore-arguments-ruling-out-siegel-zeroes/comment-page-1/\#comment-698449}{Small-prime inertness and weighted prime sums}.
\bibitem{rs} J.~Barkley Rosser and L.~Schoenfeld,
\emph{Approximate formulas for some functions of prime numbers},
Illinois J. Math. \textbf{6} (1962), 64--94.
\href{https://doi.org/10.1215/ijm/1255631807}{doi:10.1215/ijm/1255631807}.
Theorems 4, 6, 9 and 14, and Table I on printed page 90.
\bibitem{pratt} K.~Pratt, comment 698476 on \cite{tao}.
\href{https://terrytao.wordpress.com/2026/10/20/an-experimental-ai-math-challenge-using-tenth-discriminant-bounds-to-explore-arguments-ruling-out-siegel-zeroes/comment-page-1/\#comment-698476}{Norm counts, stronger divisibility, and disk interpolation}.
\bibitem{greencomment} B.~Green, comment 698463 on \cite{tao}.
\href{https://terrytao.wordpress.com/2026/10/20/an-experimental-ai-math-challenge-using-tenth-discriminant-bounds-to-explore-arguments-ruling-out-siegel-zeroes/comment-page-1/\#comment-698463}{Auxiliary-prime rank argument and Pescador attribution}.
\bibitem{kerrcomment} B.~Kerr, comment 698475 on \cite{tao}.
\href{https://terrytao.wordpress.com/2026/10/20/an-experimental-ai-math-challenge-using-tenth-discriminant-bounds-to-explore-arguments-ruling-out-siegel-zeroes/comment-page-1/\#comment-698475}{Monomial regions and higher Frobenius relations}.
See also \emph{A simpler proof that there are no Landau--Siegel zeros},
\href{https://arxiv.org/abs/2610.21099}{arXiv:2610.21099}.
\bibitem{ross} michaelmross, comment 698466 on \cite{tao}.
\href{https://terrytao.wordpress.com/2026/10/20/an-experimental-ai-math-challenge-using-tenth-discriminant-bounds-to-explore-arguments-ruling-out-siegel-zeroes/comment-page-1/\#comment-698466}{Poisson expansion of the principal quadratic form}.
\bibitem{prattsieve} K.~Pratt, comment 698460 on \cite{tao}.
\href{https://terrytao.wordpress.com/2026/10/20/an-experimental-ai-math-challenge-using-tenth-discriminant-bounds-to-explore-arguments-ruling-out-siegel-zeroes/comment-page-1/\#comment-698460}{Small-prime congruences and finite sieving}.
\bibitem{dlmf} NIST Digital Library of Mathematical Functions,
\href{https://dlmf.nist.gov/2.10}{Section 2.10 (Euler--Maclaurin)} and
\href{https://dlmf.nist.gov/5.11}{Section 5.11 (Stirling expansions
and positive-real remainder bounds)}.
\end{thebibliography}
\end{document}
